\chapter{Annales corrigées — Épreuves officielles du BEPC}

\noindent\textit{Ce chapitre reproduit quatre épreuves officielles de Mathématiques du BEPC
(Ministère des Enseignements Secondaires, Direction des Examens, des Concours et de la
Certification — MINESEC/DECC), sessions 2021, 2023, 2024 et 2025, telles qu'elles ont été
proposées aux candidats. Chaque épreuve est suivie d'un corrigé entièrement rédigé et
vérifié, pour un entraînement dans les conditions réelles de l'examen.}
\vspace{8pt}

\connecteBanner

% ═══════════════════════════════════════════════════════════════
\section{Session 2021 — Épreuve}

\noindent\textbf{Examen :} BEPC \quad\textbf{Session :} 2021 \quad\textbf{Épreuve :}
Mathématiques \quad\textbf{Durée :} 2 heures \quad\textbf{Coefficient :} 4 \quad
\textbf{Spécialités :} Toutes.
\par\smallskip\noindent
L'épreuve comporte deux parties A et B, indépendantes et obligatoires.

\subsection*{Partie A : évaluation des ressources (10 points)}

\paragraph{I. Activités numériques (5 points)}

\textbf{Exercice 1} (3 points)

\emph{A.} Pour chacune des questions suivantes, une seule des quatre réponses proposées est
juste. Reproduire et compléter le tableau ci-contre par la lettre précédant la réponse juste
pour chaque question.

\begin{enumerate}[label=\arabic*)]
\item L'écriture scientifique du nombre $A=\dfrac{21\times10^{-4}\times500\times(10^{2})^{3}}{0{,}7\times10^{8}}$
est : \\
a) $1{,}5$~; \quad b) $1{,}5\times10^{-2}$~; \quad c) $1{,}5\times10^{-1}$~; \quad
d) $1{,}5\times10^{-3}$. \hfill (0,5pt)
\item Des scientifiques ont mesuré la taille des ailes de certaines chauves-souris qui
vivent à Yaoundé.
\begin{center}
\small
\begin{tabular}{lcccccc}
\toprule
Tailles (cm) & 20 & 21 & 22 & 23 & 24 \\
\midrule
Effectifs & 22 & 20 & 28 & 17 & 13 \\
\bottomrule
\end{tabular}
\end{center}
La moyenne de cette série statistique est : \\
a) $21{,}79$~; \quad b) $22{,}79$~; \quad c) $21{,}79$~; \quad d) $23{,}79$. \hfill (0,5pt)
\item La forme factorisée de $B=(5x-9)^{2}-(5x-9)$ est : \\
a) $(5x-9)(5x+8)$~; \quad b) $(5x-9)(5x-9)$~; \quad c) $5(5x-9)(x-1)$~; \quad
d) $5(5x-9)(x-2)$. \hfill (0,5pt)
\item On donne $A=\dfrac34+\dfrac52\div\dfrac56-2\left(\dfrac32\right)^{2}$. L'écriture
irréductible de $A$ est : \\
a) $\dfrac34$~; \quad b) $-\dfrac54$~; \quad c) $-\dfrac34$~; \quad d) $\dfrac54$. \hfill
(0,5pt)
\end{enumerate}

\emph{B.} Pour chacune des affirmations ci-après, répondre par Vrai (V) ou Faux (F).
\begin{enumerate}[label=\arabic*)]
\item Le nombre $C=(3\sqrt2+2)(3\sqrt2-2)$ est un entier naturel. \hfill (0,5pt)
\item Si $2{,}23<\sqrt5<2{,}24$ alors $0{,}76<3-\sqrt5<0{,}77$. \hfill (0,5pt)
\end{enumerate}

\textbf{Exercice 2} (2 points)
\begin{enumerate}[label=\arabic*)]
\item En utilisant l'algorithme d'Euclide, montrer que le PGCD de $378$ et $270$ est $54$.
Faire apparaître sur votre copie les calculs intermédiaires. \hfill (1pt)
\item Pour une kermesse, un comité des fêtes dispose de $378$ cannettes de jus et de $270$
cannettes de bière. Il veut faire le plus grand nombre de lots identiques en utilisant
toutes les cannettes.
\begin{enumerate}[label=\alph*)]
\item Combien de ces lots identiques pourra-t-il placer sur un certain nombre de tables ?
\hfill (0,5pt)
\item Quelle sera la composition de chacun de ces lots ? \hfill (0,5pt)
\end{enumerate}
\end{enumerate}

\paragraph{II. Activités géométriques (5 points)}

\textbf{Exercice 1} (3 points) \\
L'unité de longueur est le centimètre. $ABC$ est un triangle tel que : $AB=4{,}5$~;
$AC=7{,}5$~; $BC=6$.
\begin{enumerate}[label=\arabic*)]
\item Tracer le triangle $ABC$ sur votre copie et démontrer que ce triangle est rectangle
en $B$. \hfill (0,75pt)
\item $E$ est un point du segment $[AB]$ tel que $AE=1{,}5$. $F$ est un point du segment
$[AC]$ tel que $AF=2{,}5$.
\begin{enumerate}[label=\alph*)]
\item Placer $E$ et $F$ sur la figure. \hfill (0,5pt)
\item Démontrer que les droites $(EF)$ et $(BC)$ sont parallèles. \hfill (0,75pt)
\end{enumerate}
\item Calculer $EF$, puis en déduire l'aire du trapèze $EBCF$. \hfill (1pt)
\end{enumerate}

\textbf{Exercice 2} (2 points) \\
On considère trois droites $(D_1)$, $(D_2)$ et $(D_3)$ construites dans le repère
orthonormé ci-dessous. On donne trois équations cartésiennes : $2x-y-2=0$~; $x+2y+2=0$ et
$2x-y+2=0$.

\begin{illus}
\begin{tikzpicture}
\begin{axis}[width=7.5cm,height=7.5cm,axis lines=middle,xlabel={\small $x$},ylabel={\small $y$},
  xmin=-4,xmax=5,ymin=-5,ymax=5,xtick={-4,-3,-2,-1,1,2,3,4,5},ytick={-5,-4,-3,-2,-1,1,2,3,4,5},
  grid=both,minor grid style={gray!12},major grid style={gray!25},
  tick label style={font=\scriptsize}]
\addplot[domain=-4:5,onbo,thick,samples=2]{2*x+2};
\addplot[domain=-4:5,onbgreen,thick,samples=2]{2*x-2};
\addplot[domain=-4:5,onbink,thick,samples=2]{-0.5*x-1};
\node[onbo] at (axis cs:1.2,4.55) {\scriptsize $D_1$};
\node[onbgreen] at (axis cs:3.6,4.55) {\scriptsize $D_2$};
\node[onbink] at (axis cs:-3.4,0.85) {\scriptsize $D_3$};
\end{axis}
\end{tikzpicture}
\fig{Figure — Droites $(D_1)$, $(D_2)$, $(D_3)$ dans le repère orthonormé.}
\end{illus}

\begin{enumerate}[label=\arabic*)]
\item Associer à chacune des droites, son équation cartésienne. \hfill (0,75pt)
\item Déterminer les coefficients directeurs des droites $(D_1)$ et $(D_2)$ et en déduire
qu'elles sont parallèles. \hfill (0,5pt)
\item[3)] a) Déterminer le couple $(x\,;y)$ de nombres réels vérifiant :
$\begin{cases}2x-y=-2\\x+2y=-2\end{cases}$. \hfill (0,5pt) \\
b) En déduire une interprétation graphique de ce résultat. \hfill (0,25pt)
\end{enumerate}

\subsection*{Partie B : évaluation des compétences (10 points)}

\textbf{Situation.} Suite à une baisse des salaires dans l'entreprise où il est employé,
Ateba décide de réduire les consommations d'eau, d'électricité et de gaz de sa famille.

Chaque mois, la famille Ateba consomme $80$ m$^{3}$ d'eau de la société « EAU-POUR-TOUS »
et $385$ kWh d'électricité de la société « SONELEC ». À la fin de chaque mois, Ateba paie
une facture d'eau et une facture d'électricité. Chacune des deux factures comporte le coût
de la consommation, augmenté d'une T.V.A. représentant $19{,}25\,\%$ du coût de cette
consommation. $1$ m$^{3}$ d'eau coûte $365$ FCFA et $1$ kWh d'électricité coûte $65$ FCFA.

Pour diminuer la consommation d'eau, Ateba réaménage un puits dans sa concession. Pour
réduire sa consommation d'électricité, Ateba remplace les ampoules incandescentes par des
ampoules économiques.

La quantité de gaz utilisée par mois est en moyenne de $24$ litres, à raison de $550$ FCFA
le litre. Pour diminuer cette consommation, Mme Ateba prépare certains mets au feu de bois.
Une étude a montré que $4$ litres de gaz utilisé équivalent à un fagot de bois standard du
marché local.

\textbf{Tâches :}
\begin{enumerate}[label=\arabic*)]
\item Déterminer le volume minimal d'eau en m$^{3}$ provenant du puits que la famille Ateba
doit utiliser pour payer une facture d'eau d'un montant inférieur ou égal à $25\,000$ FCFA.
\hfill (3pts)
\item Déterminer la consommation minimale des ampoules économiques (en kWh), pour que la
famille Ateba paie une facture d'électricité d'un montant inférieur ou égal à $26\,000$
FCFA. \hfill (3pts)
\item Déterminer le nombre minimal de fagots de bois que la famille Ateba doit utiliser
pour payer une facture de gaz d'un montant inférieur ou égal à $7\,000$ FCFA. \hfill (3pts)
\end{enumerate}
\noindent N.B. On donnera les résultats arrondis à l'unité. \hfill \textbf{Présentation : 1pt}

% ═══════════════════════════════════════════════════════════════
\section{Session 2021 — Corrigé}
\begin{multicols}{2}

\corrige{— Partie A, Activités numériques, Ex. 1.A.1}
$A=\dfrac{21\times10^{-4}\times500\times10^{6}}{0{,}7\times10^{8}}=\dfrac{21\times500}{0{,}7}\times10^{-4+6-8}=15\,000\times10^{-6}=1{,}5\times10^{-2}$.
Réponse : \textbf{b}.

\corrige{— Ex. 1.A.2}
Effectif total $=22+20+28+17+13=100$.\\
Moyenne $=\dfrac{440+420+616+391+312}{100}$\\
$=\dfrac{2\,179}{100}=21{,}79$.
Réponses \textbf{a} et \textbf{c} (identiques dans l'énoncé) sont toutes deux correctes.

\corrige{— Ex. 1.A.3}
$B=(5x-9)[(5x-9)-1]=(5x-9)(5x-10)=5(5x-9)(x-2)$. Réponse : \textbf{d}.

\corrige{— Ex. 1.A.4}
$A=\dfrac34+\dfrac52\times\dfrac65-2\times\dfrac94=\dfrac34+3-\dfrac92=\dfrac34+\dfrac{12}4-\dfrac{18}4=-\dfrac34$.
Réponse : \textbf{c}.

\corrige{— Ex. 1.B.1}
$C=(3\sqrt2)^{2}-2^{2}=18-4=14$ : entier naturel. \textbf{Vrai}.

\corrige{— Ex. 1.B.2}
$2{,}23<\sqrt5<2{,}24\Rightarrow-2{,}24<-\sqrt5<-2{,}23\Rightarrow0{,}76<3-\sqrt5<0{,}77$.
\textbf{Vrai}.

\corrige{— Ex. 2.1}
$378=1\times270+108$~; $270=2\times108+54$~; $108=2\times54+0$. Le dernier reste non nul est
$54$ : $\text{PGCD}(378\,;270)=54$.

\corrige{— Ex. 2.2}
a) Nombre de lots $=54$. \\
b) Chaque lot contient $378\div54=7$ cannettes de jus et $270\div54=5$ cannettes de bière.

\corrige{— Partie A, Activités géométriques, Ex. 1.1}
$AC^{2}=7{,}5^{2}=56{,}25$ et $AB^{2}+BC^{2}=4{,}5^{2}+6^{2}=20{,}25+36=56{,}25$. Égalité :
le triangle $ABC$ est rectangle en $B$ (réciproque de Pythagore).

\corrige{— Ex. 1.2}
b) $\dfrac{AE}{AB}=\dfrac{1{,}5}{4{,}5}=\dfrac13$ et $\dfrac{AF}{AC}=\dfrac{2{,}5}{7{,}5}=\dfrac13$ :
rapports égaux, donc par la réciproque de Thalès, $(EF)\parallel(BC)$.

\corrige{— Ex. 1.3}
$EF=BC\times\dfrac13=6\times\dfrac13=2$ cm. Aire du triangle $ABC=\dfrac12\times AB\times BC=\dfrac12\times4{,}5\times6=13{,}5$ cm$^{2}$.
Aire du triangle $AEF=13{,}5\times\left(\dfrac13\right)^{2}=1{,}5$ cm$^{2}$. Aire du trapèze
$EBCF=13{,}5-1{,}5=12$ cm$^{2}$.

\corrige{— Ex. 2.1}
$(D_1)$ passe par $(0\,;2)$ et $(-1\,;0)$ : équation $2x-y+2=0$. $(D_2)$ passe par
$(0\,;-2)$ et $(1\,;0)$ : équation $2x-y-2=0$. $(D_3)$ passe par $(0\,;-1)$ et $(-2\,;0)$ :
équation $x+2y+2=0$.

\corrige{— Ex. 2.2}
$(D_1)$ : $y=2x+2$, coefficient directeur $2$. $(D_2)$ : $y=2x-2$, coefficient directeur
$2$. Coefficients directeurs égaux : $(D_1)\parallel(D_2)$.

\corrige{— Ex. 2.3}
a) $2x-y=-2\Rightarrow y=2x+2$. En substituant : $x+2(2x+2)=-2\Rightarrow5x+4=-2\Rightarrow x=-\dfrac65$,
puis $y=2\times\left(-\dfrac65\right)+2=-\dfrac25$. Couple solution : $\left(-\dfrac65\,;-\dfrac25\right)$. \\
b) Les deux équations du système sont celles de $(D_1)$ et $(D_3)$ : le couple solution est
donc les coordonnées du point d'intersection de $(D_1)$ et $(D_3)$.

\corrige{— Partie B, Tâche 1}
Facture d'eau $=365\times V_{\text{réseau}}\times1{,}1925\le25\,000\Rightarrow
V_{\text{réseau}}\le\dfrac{25\,000}{365\times1{,}1925}\approx57{,}44$ m$^{3}$.
Volume minimal du puits $=80-57{,}44\approx22{,}56$, soit \textbf{$23$ m$^{3}$} (arrondi à
l'unité).

\corrige{— Tâche 2}
Facture d'électricité $=65\times C_{\text{payante}}\times1{,}1925\le26\,000\Rightarrow
C_{\text{payante}}\le\dfrac{26\,000}{65\times1{,}1925}\approx335{,}43$ kWh. Consommation
minimale des ampoules économiques $=385-335{,}43\approx49{,}57$, soit \textbf{$50$ kWh}
(arrondi à l'unité).

\corrige{— Tâche 3}
Facture de gaz $=550\times V_{\text{gaz restant}}\le7\,000\Rightarrow
V_{\text{gaz restant}}\le\dfrac{7\,000}{550}\approx12{,}73$ L. Volume remplacé par du bois
$=24-12{,}73\approx11{,}27$ L, soit $\dfrac{11{,}27}{4}\approx2{,}82$ fagots, donc
\textbf{$3$ fagots} minimum (arrondi à l'unité supérieure pour garantir la facture visée).

\end{multicols}

% ═══════════════════════════════════════════════════════════════
\section{Session 2023 — Épreuve}

\noindent\textbf{Examen :} BEPC \quad\textbf{Session :} 2023 \quad\textbf{Épreuve :}
Mathématiques \quad\textbf{Durée :} 2 heures \quad\textbf{Coefficient :} 4 \quad
\textbf{Série :} Toutes.

\subsection*{Partie A : évaluation des ressources (10 points)}

\paragraph{Activités numériques (5 points)}

\textbf{Exercice 1} (2 points)
\begin{enumerate}[label=\arabic*)]
\item Calcule le nombre $A=\dfrac{25}{18}-\dfrac79\times\left(-\dfrac5{14}+\dfrac8{21}\right)$
et donne le résultat sous forme d'une fraction irréductible. \hfill (1pt)
\item Soit $B=(3-2\sqrt3)^{2}$. Montre que $B=21-12\sqrt3$. \hfill (0,5pt)
\item Choisis la bonne réponse parmi les quatre qui sont proposées : \\
a) $3-2\sqrt3$~; \quad b) $3+2\sqrt3$~; \quad c) $-3+2\sqrt3$~; \quad d) $-3-2\sqrt3$.
\hfill (0,5pt)
\end{enumerate}

\textbf{Exercice 2} (1,5 points)
\begin{enumerate}[label=\arabic*)]
\item On considère l'expression $C=(x-1)^{2}+(x-1)(2x-3)$. Mets $C$ sous la forme d'un
produit de facteurs du premier degré. \hfill (0,75pt)
\item On considère la fraction rationnelle $Q=\dfrac{(x-1)(3x-4)}{(x-1)(x+2)}$. Donne la
condition d'existence d'une valeur numérique de $Q$ puis simplifie $Q$. \hfill (0,75pt)
\end{enumerate}

\textbf{Exercice 3} (1,5 points) \\
L'histogramme ci-dessous représente la série statistique des âges regroupés en classes,
des $30$ enfants d'un club de sport.

\begin{illus}
\begin{tikzpicture}[scale=1]
\draw[onbmuted,->] (-0.3,0) -- (6.5,0);
\draw[onbmuted,->] (0,-0.3) -- (0,7.8);
\node[right,font=\scriptsize] at (6.5,0) {Âge (années)};
\node[above,font=\scriptsize] at (0,7.8) {Effectifs};
\draw[fill=onbo!55,draw=onbo,line width=0.8pt] (0,0) rectangle (1,2);
\draw[fill=onbo!55,draw=onbo,line width=0.8pt] (1,0) rectangle (2,3);
\draw[fill=onbo!55,draw=onbo,line width=0.8pt] (2,0) rectangle (3,7);
\draw[fill=onbo!55,draw=onbo,line width=0.8pt] (3,0) rectangle (4,6);
\draw[fill=onbo!55,draw=onbo,line width=0.8pt] (4,0) rectangle (5,7);
\draw[fill=onbo!55,draw=onbo,line width=0.8pt] (5,0) rectangle (6,5);
\foreach \x/\lbl in {0/8,1/9,2/10,3/11,4/12,5/13,6/14}
  \node[below,font=\scriptsize] at (\x,0) {\lbl};
\node[above,font=\scriptsize] at (0.5,2) {2};
\node[above,font=\scriptsize] at (1.5,3) {3};
\node[above,font=\scriptsize] at (2.5,7) {7};
\node[above,font=\scriptsize] at (3.5,6) {6};
\node[above,font=\scriptsize] at (4.5,7) {7};
\node[above,font=\scriptsize] at (5.5,5) {5};
\end{tikzpicture}
\fig{Figure — Histogramme des âges des 30 enfants du club de sport.}
\end{illus}

\begin{enumerate}[label=\arabic*)]
\item Détermine l'amplitude des classes de cette série. \hfill (0,25pt)
\item Donne une classe modale de cette série statistique. \hfill (0,25pt)
\item Calcule la moyenne des âges des enfants de ce club. \hfill (1pt)
\end{enumerate}

\paragraph{Activités géométriques (5 points)}

\textbf{Exercice 1} (1,75 points) \\
Le plan est muni d'un repère orthonormé $(O,I,J)$. On considère les points $A$ et $B$ de
coordonnées respectives $(3\,;2)$~; $(-2\,;1)$.
\begin{enumerate}[label=\arabic*)]
\item Calcule les coordonnées du vecteur $\overrightarrow{AB}$. \hfill (0,5pt)
\item Détermine une équation cartésienne de la droite $(AB)$. \hfill (0,75pt)
\item Réponds par VRAI ou FAUX. Les droites $(D)$ et $(L)$ d'équations respectives
$x-5y+7=0$ et $5x+y+9=0$ sont perpendiculaires. \hfill (0,5pt)
\end{enumerate}

\textbf{Exercice 2} (1,75 points) \\
Sur la figure ci-dessous, le triangle $EFG$, rectangle en $E$, est tel que $EF=6$ cm~;
$EG=8$ cm. $M$ et $N$ sont respectivement des points de $[EF]$ et $[EG]$ tels que
$EM=1{,}5$ cm et $EN=2$ cm.

\begin{illus}
\begin{tikzpicture}[scale=0.7]
\draw[onbline,line width=1.3pt] (0,0) -- (6,0) -- (0,8) -- cycle;
\draw[onbmuted,line width=0.7pt] (0,0.4)--(0.4,0.4)--(0.4,0);
\draw[onbo,line width=1.1pt] (0,1.5) -- (2,1.5);
\filldraw[onbo] (0,1.5) circle(1.6pt) node[left,font=\scriptsize]{$M$};
\filldraw[onbo] (2,1.5) circle(1.6pt) node[right,font=\scriptsize]{$N$};
\node[below left,font=\scriptsize] at (0,0) {$E$};
\node[below,font=\scriptsize] at (6,0) {$F$};
\node[above,font=\scriptsize] at (0,8) {$G$};
\end{tikzpicture}
\fig{Figure — Triangle $EFG$ rectangle en $E$, avec $M\in[EF]$, $N\in[EG]$.}
\end{illus}

\begin{enumerate}[label=\arabic*)]
\item Montre que les droites $(FG)$ et $(MN)$ sont parallèles. \hfill (0,75pt)
\item Calcule $\tan\widehat{EGF}$ et déduis-en l'arrondi à $1^{\circ}$ près de la mesure de
l'angle $\widehat{EGF}$. \hfill (0,5pt)
\item Le triangle $EFG$ est l'image du triangle $EMN$ par une homothétie. Donne le centre
et le rapport de cette homothétie. \hfill (0,5pt)
\end{enumerate}

\textbf{Exercice 3} (1,5 points) \\
La figure ci-dessous représente un grand cône de révolution coupé par un plan parallèle à
la base et passant par $O'$ et $B$. On désigne par $v$ et $a$ le volume du petit cône et
l'aire du disque de rayon $[O'B]$. On donne $v=10$ cm$^{3}$~; $a=5$ cm$^{2}$ et
$\dfrac{SO'}{SO}=\dfrac13$.

\begin{illus}
\begin{tikzpicture}[scale=0.7]
\coordinate (S) at (0,4.2);
\coordinate (Ob) at (0,0);
\coordinate (Oc) at (0,1.4);
\draw[onbline,line width=0.8pt] (2.0,0) arc (0:-180:2.0 and 0.55);
\draw[onbline,dashed,line width=0.6pt] (2.0,0) arc (0:180:2.0 and 0.55);
\draw[onbline,line width=1.1pt] (S) -- (2.0,0) node[below,font=\scriptsize]{$A$};
\draw[onbline,line width=1.1pt] (S) -- (-2.0,0);
\draw[onbo,line width=1pt,fill=onbsoft] (1.333,1.4) arc (0:360:1.333 and 0.37);
\node[onbo,right,font=\scriptsize] at (1.333,1.4) {$B$};
\node[left,font=\scriptsize] at (Oc) {$O'$};
\node[left,font=\scriptsize] at (Ob) {$O$};
\node[above,font=\scriptsize] at (S) {$S$};
\end{tikzpicture}
\fig{Figure — Grand cône de sommet $S$, coupé en $O'B$ parallèlement à la base $OA$.}
\end{illus}

\begin{enumerate}[label=\arabic*)]
\item Calcule le volume $\mathcal{V}$ du grand cône. \hfill (0,75pt)
\item Calcule l'aire $\mathcal{A}$ de la base du grand cône. \hfill (0,75pt)
\end{enumerate}

\subsection*{Partie B : évaluation des compétences (10 points)}

\textbf{Situation.} ANGO a un terrain rectangulaire de dimensions $546$ mètres et $510$
mètres, qu'il voudrait clôturer pour besoin de sécurité. Afin de poser un grillage, il doit
planter des poteaux régulièrement espacés. Par souci de limiter les dépenses, il souhaite
utiliser le plus petit nombre possible de poteaux et que la distance entre deux poteaux
soit un nombre entier de mètres. Il place un poteau à chaque coin du terrain. ANGO commence
les travaux par l'une des longueurs du terrain.

Pour l'achat des poteaux, il s'adresse à deux vendeurs ALPHA et BETA qui lui font les
offres suivantes :
\begin{itemize}
\item Vendeur ALPHA : un poteau à $1\,400$ FCFA et $20\,000$ FCFA pour les frais de
livraison ;
\item Vendeur BETA : un poteau à $1\,600$ FCFA et la livraison gratuite.
\end{itemize}
Pour opérer son choix, ANGO voudrait savoir le nombre minimal de poteaux pour lequel
l'offre du vendeur ALPHA est la plus avantageuse des deux.

Pour l'exploitation des $278\,460$ m$^{2}$ de son terrain, il a prévu trois parcelles de
culture : la première pour le palmier à huile, la deuxième représentant les $\dfrac45$ de
la première pour le maïs, et les $54\,360$ m$^{2}$ restants pour les arbres fruitiers. Pour
un meilleur rendement, les techniciens lui indiquent qu'il faut réserver $35$ m$^{2}$ par
palmier, mais il a oublié la superficie de la parcelle réservée à la culture du palmier à
huile.

\textbf{Tâches :}
\begin{enumerate}[label=\arabic*)]
\item Calcule le nombre minimal de poteaux que ANGO doit utiliser pour l'une des longueurs
de son terrain. \hfill (3pts)
\item Calcule le nombre minimal de poteaux pour lequel l'offre du vendeur ALPHA est plus
avantageuse que celle du vendeur BETA. \hfill (3pts)
\item Calcule le nombre de palmiers à huile que ANGO peut planter sur la première parcelle.
\hfill (3pts)
\end{enumerate}
\noindent\hfill\textbf{Présentation générale : 1 point}

% ═══════════════════════════════════════════════════════════════
\section{Session 2023 — Corrigé}
\begin{multicols}{2}

\corrige{— Activités numériques, Ex. 1.1}
$-\dfrac5{14}+\dfrac8{21}=\dfrac{-15+16}{42}=\dfrac1{42}$ ; $\dfrac79\times\dfrac1{42}=\dfrac1{54}$ ;
$A=\dfrac{25}{18}-\dfrac1{54}=\dfrac{75-1}{54}=\dfrac{74}{54}=\dfrac{37}{27}$.

\corrige{— Ex. 1.2}
$B=(3-2\sqrt3)^{2}=9-12\sqrt3+4\times3=9-12\sqrt3+12=21-12\sqrt3$.

\corrige{— Ex. 1.3}
$\sqrt B=|3-2\sqrt3|$. Comme $2\sqrt3\approx3{,}46>3$, on a $3-2\sqrt3<0$, donc
$\sqrt B=-(3-2\sqrt3)=-3+2\sqrt3$. Réponse : \textbf{c}.

\corrige{— Ex. 2.1}
$C=(x-1)[(x-1)+(2x-3)]=(x-1)(3x-4)$.

\corrige{— Ex. 2.2}
Condition d'existence : $x\ne1$ et $x\ne-2$. $Q=\dfrac{3x-4}{x+2}$.

\corrige{— Ex. 3.1}
Amplitude des classes $=1$ an.

\corrige{— Ex. 3.2}
Effectif maximal $=7$, atteint par les classes $[10\,;11[$ et $[12\,;13[$ : la série est
bimodale (deux classes modales).

\corrige{— Ex. 3.3}
Avec les centres de classes $8{,}5$, $9{,}5$, $10{,}5$, $11{,}5$, $12{,}5$, $13{,}5$ :\\
moyenne $=\dfrac{17+28{,}5+73{,}5+69+87{,}5+67{,}5}{30}$\\
$=\dfrac{343}{30}\approx11{,}43$ ans.

\corrige{— Activités géométriques, Ex. 1.1}
$\overrightarrow{AB}=(-2-3\,;1-2)=(-5\,;-1)$.

\corrige{— Ex. 1.2}
Droite de direction $(-5\,;-1)$ passant par $A(3\,;2)$ :
$\dfrac{y-2}{x-3}=\dfrac{-1}{-5}=\dfrac15\Rightarrow5(y-2)=x-3\Rightarrow x-5y+7=0$.

\corrige{— Ex. 1.3}
$(D)$ : $x-5y+7=0$, coefficient directeur $\dfrac15$. $(L)$ : $5x+y+9=0$, coefficient
directeur $-5$. Produit $=\dfrac15\times(-5)=-1$ : les droites sont perpendiculaires.
\textbf{Vrai}.

\corrige{— Ex. 2.1}
$\dfrac{EM}{EF}=\dfrac{1{,}5}6=0{,}25$ et $\dfrac{EN}{EG}=\dfrac28=0{,}25$ : rapports
égaux, donc par la réciproque de Thalès, $(MN)\parallel(FG)$.

\corrige{— Ex. 2.2}
$\tan\widehat{EGF}=\dfrac{EF}{EG}=\dfrac68=0{,}75$, donc
$\widehat{EGF}=\arctan(0{,}75)\approx37^{\circ}$.

\corrige{— Ex. 2.3}
L'homothétie de centre $E$ et de rapport $k=\dfrac{EF}{EM}=\dfrac6{1{,}5}=4$ transforme $M$
en $F$ et $N$ en $G$ : centre $E$, rapport $4$.

\corrige{— Ex. 3.1}
Rapport de réduction $k=\dfrac{SO'}{SO}=\dfrac13$, donc $\dfrac{v}{\mathcal V}=k^{3}=\dfrac1{27}$ :
$\mathcal V=27\times v=27\times10=270$ cm$^{3}$.

\corrige{— Ex. 3.2}
$\dfrac{a}{\mathcal A}=k^{2}=\dfrac19$, donc $\mathcal A=9\times a=9\times5=45$ cm$^{2}$.

\corrige{— Partie B, Tâche 1}
$\text{PGCD}(546\,;510)$ : $546=1\times510+36$~; $510=14\times36+6$~; $36=6\times6+0$ :
$\text{PGCD}=6$. Espacement maximal $=6$ m. Nombre de poteaux sur une longueur de $546$ m :
$546\div6+1=92$ poteaux.

\corrige{— Tâche 2}
Coût ALPHA $=1\,400n+20\,000$ ; coût BETA $=1\,600n$. ALPHA avantageux si
$1\,400n+20\,000<1\,600n\Leftrightarrow20\,000<200n\Leftrightarrow n>100$. Nombre minimal :
$n=101$ poteaux.

\corrige{— Tâche 3}
Aire du terrain $=546\times510=278\,460$ m$^{2}$ (vérifiée). Soit $P$ l'aire du palmier :
$P+\dfrac45P+54\,360=278\,460\Rightarrow\dfrac95P=224\,100\Rightarrow P=124\,500$ m$^{2}$.
Nombre de palmiers $=124\,500\div35\approx3\,557{,}14$ : on ne peut planter que
$\textbf{3\,557}$ palmiers (il reste $5$ m$^{2}$ inutilisés).

\end{multicols}

% ═══════════════════════════════════════════════════════════════
\section{Session 2024 — Épreuve}

\noindent\textbf{Examen :} BEPC \quad\textbf{Session :} 2024 \quad\textbf{Épreuve :}
Mathématiques \quad\textbf{Durée :} 2 heures \quad\textbf{Coefficient :} 4 \quad
\textbf{Série :} Toutes.

\subsection*{Partie A : évaluation des ressources (10 points)}

\paragraph{Activités numériques (5 points)}

\textbf{Exercice 1} (1,5 points) \\
On considère le nombre $A=\dfrac{2+\sqrt3}{2-\sqrt3}$.
\begin{enumerate}[label=\arabic*)]
\item Calculer $(2+\sqrt3)^{2}$. \hfill (0,5pt)
\item Montrer que $A=7+4\sqrt3$. \hfill (0,5pt)
\item Sachant que $1{,}732<\sqrt3<1{,}733$, déterminer un encadrement de $A$ par deux
nombres décimaux d'ordre $2$. \hfill (0,5pt)
\end{enumerate}

\textbf{Exercice 2} (1,5 points)
\begin{enumerate}[label=\arabic*)]
\item On considère l'expression $B=(x-2)^{2}+(x-2)(x+4)$. Écrire $B$ sous la forme d'un
produit de facteurs du premier degré. \hfill (0,75pt)
\item On considère la fraction rationnelle $C=\dfrac{(x-2)(2x+1)}{(x-2)(x+3)}$. Donner la
condition d'existence d'une valeur numérique de $C$ puis simplifier $C$. \hfill (0,75pt)
\end{enumerate}

\textbf{Exercice 3} (2 points) \\
Le tableau suivant récapitule les notes sur $20$ en mathématiques de $50$ élèves d'une
classe de troisième, avec une donnée manquante.
\begin{center}
\small
\begin{tabular}{lccccc}
\toprule
Notes & $[0\,;4[$ & $[4\,;8[$ & $[8\,;12[$ & $[12\,;16[$ & $[16\,;20[$ \\
\midrule
Effectifs & 13 & 15 & 10 & \ ? & 5 \\
\bottomrule
\end{tabular}
\end{center}
\begin{enumerate}[label=\arabic*)]
\item Quel est le caractère étudié et quelle est sa nature ? \hfill (0,5pt)
\item Justifier que l'effectif de la classe $[12\,;16[$ est $7$. \hfill (0,25pt)
\item Quelle est la classe modale de cette série statistique ? \hfill (0,25pt)
\item Calculer la moyenne des notes sur $20$ en mathématiques de cette classe. \hfill (1pt)
\end{enumerate}

\paragraph{Activités géométriques (5 points)}

\textbf{Exercice 1} (1,5 point) \\
Le plan est muni d'un repère orthonormé $(O,I,J)$. On considère les points $A$, $B$ et $C$
de coordonnées respectives $(2\,;2)$~; $(-3\,;1)$ et $(4\,;-2)$. Répondre par Vrai ou Faux à
chacune des questions suivantes :
\begin{enumerate}[label=\arabic*)]
\item Le couple de coordonnées du vecteur $\overrightarrow{AB}$ est : $(-5\,;-1)$. \hfill
(0,5pt)
\item La distance $AC$ est égale à $3\sqrt5$. \hfill (0,5pt)
\item Une équation cartésienne de la droite $(AB)$ est : $x-5y+8=0$. \hfill (0,5pt)
\end{enumerate}

\textbf{Exercice 2} (2 points) \\
Sur la figure ci-dessous, $ABC$ est un triangle rectangle en $C$ tel que : $AB=5$ cm~;
$BC=3$ cm et $AC=4$ cm. $E$ est le point du segment $[AB]$ tel que $AE=4$ cm. La droite
passant par $E$ et parallèle à $(BC)$ coupe le segment $[AC]$ en $F$. $(\mathcal C)$ est le
cercle de centre $O$, circonscrit au triangle $ABC$. $M$ est un point du cercle
$(\mathcal C)$.

\begin{illus}
\begin{tikzpicture}[scale=0.85]
\draw[onbline,line width=0.7pt] (0,0) circle (2.1cm);
\coordinate (A) at (2.1,0);
\coordinate (B) at (-1.05,-1.818);
\coordinate (C) at (-0.85,1.94);
\coordinate (M) at (0.7,-1.98);
\draw[onbo,line width=1.2pt] (A) -- (B) -- (C) -- cycle;
\draw[onbmuted,line width=0.6pt] (A) -- (M) -- (B);
\coordinate (E) at ($(A)!0.8!(B)$);
\coordinate (F) at ($(A)!0.8!(C)$);
\draw[onbgreen,line width=1pt] (E) -- (F);
\filldraw[onbo] (A) circle(1.6pt) node[right,font=\scriptsize]{$A$};
\filldraw[onbo] (B) circle(1.6pt) node[below left,font=\scriptsize]{$B$};
\filldraw[onbo] (C) circle(1.6pt) node[above left,font=\scriptsize]{$C$};
\filldraw[onbgreen] (E) circle(1.4pt) node[below,font=\scriptsize]{$E$};
\filldraw[onbgreen] (F) circle(1.4pt) node[left,font=\scriptsize]{$F$};
\filldraw[onbmuted] (M) circle(1.4pt) node[below right,font=\scriptsize]{$M$};
\filldraw[black] (0,0) circle(1.2pt) node[above right,font=\scriptsize]{$O$};
\end{tikzpicture}
\fig{Figure — Triangle $ABC$ inscrit dans le cercle $(\mathcal C)$, avec $E\in[AB]$, $F\in[AC]$, $EF\parallel BC$.}
\end{illus}

\begin{enumerate}[label=\arabic*)]
\item Calculer $\cos\widehat{BAC}$ et en déduire l'arrondi à $1^{\circ}$ près de la mesure
de l'angle $\widehat{BAC}$. \hfill (0,75pt)
\item Justifier que les angles $\widehat{BAC}$ et $\widehat{BMC}$ ont la même mesure. \hfill
(0,5pt)
\item Calculer la distance $EF$. \hfill (0,75pt)
\end{enumerate}

\textbf{Exercice 3} (1,5 point) \\
La figure ci-dessous représente une pyramide régulière $SABCD$, de hauteur $SM=6$ cm et
dont la base est carrée de côté $AB=4$ cm. On coupe cette pyramide suivant un plan
parallèle à la base et passant par le point $N$, tel que $SN=\dfrac12SM$.

\begin{illus}
\begin{tikzpicture}[scale=0.9]
\coordinate (O) at (0,0);
\coordinate (F) at (0,-1.0);
\coordinate (K) at (0,1.0);
\coordinate (R) at (1.9,-0.35);
\coordinate (L) at (-1.9,0.35);
\coordinate (S) at (0,3.0);
\coordinate (N) at (0,1.5);
\coordinate (Fr) at ($(F)!0.5!(R)$);
\coordinate (Fl) at ($(F)!0.5!(L)$);
\coordinate (Kr) at ($(K)!0.5!(R)$);
\coordinate (Kl) at ($(K)!0.5!(L)$);
\draw[onbline,dashed,line width=0.7pt] (R) -- (K) -- (L);
\draw[onbline,line width=1pt] (F) -- (R);
\draw[onbline,line width=1pt] (L) -- (F);
\draw[onbline,dashed,line width=0.7pt] (S) -- (K);
\draw[onbo,line width=1.2pt] (S) -- (F);
\draw[onbo,line width=1.2pt] (S) -- (R);
\draw[onbo,line width=1.2pt] (S) -- (L);
\draw[onbgreen,dashed,line width=0.8pt] (Kr) -- (Fr) -- (Fl) -- (Kl) -- cycle;
\node[onbgreen,right,font=\scriptsize] at (Fr) {$G$};
\node[onbgreen,left,font=\scriptsize] at (Fl) {$H$};
\node[onbgreen,right,font=\scriptsize] at (Kr) {$F$};
\node[onbgreen,left,font=\scriptsize] at (Kl) {$E$};
\node[onbmuted,left,font=\scriptsize] at (N) {$N$};
\filldraw[onbmuted] (N) circle(1.3pt);
\node[below,font=\scriptsize] at (F) {$M$};
\node[above,font=\scriptsize] at (S) {$S$};
\end{tikzpicture}
\fig{Figure — Pyramide $SABCD$ coupée en $N$ (milieu de $[SM]$) : petite pyramide $SEFGH$.}
\end{illus}

\begin{enumerate}[label=\arabic*)]
\item Calculer le volume $V$ de la pyramide $SABCD$. \hfill (0,75pt)
\item En déduire le volume $v$ de la pyramide réduite $SEFGH$. \hfill (0,75pt)
\end{enumerate}

\subsection*{Partie B : évaluation des compétences (10 points)}

\textbf{Situation.} Le propriétaire d'un hôtel voudrait effectuer des travaux de
réfection. Son comptable a envisagé de faire financer les $\dfrac25$ du montant total du
budget des travaux par une banque, le $\dfrac14$ du montant total de ce budget par une
partie des recettes, et le reste, soit la somme de $1\,225\,000$ FCFA, par ses économies
personnelles. Avant de commencer les travaux, il voudrait connaître le montant exact du
budget total des travaux, mais son comptable est absent.

Une partie des travaux concerne le renouvellement du carrelage d'une cuisine de forme
rectangulaire de longueur $4{,}80$ mètres et de largeur $3$ mètres. Il envisage utiliser des
carreaux de forme carrée ayant la plus grande dimension possible, que l'on posera sans
espaces. Il voudrait connaître le nombre minimal de carreaux à acheter.

Une autre partie des travaux concerne la peinture des murs. Il a pour cela acheté un total
de $10$ seaux de peinture comprenant des seaux de peinture à huile et des seaux de peinture
à eau. La quincaillerie lui a vendu un seau de peinture à eau à $40\,000$ FCFA et un seau de
peinture à huile à $60\,000$ FCFA, pour un montant total de $440\,000$ FCFA. Au moment de la
livraison, il voudrait vérifier le nombre de seaux de peinture de chaque type.

\textbf{Tâches :}
\begin{enumerate}[label=\arabic*)]
\item Calculer le montant total du budget des travaux de réfection. \hfill (3pts)
\item Calculer le nombre minimal de carreaux à acheter pour le renouvellement du carrelage.
\hfill (3pts)
\item Calculer le nombre de seaux de peinture de chaque type qui ont été achetés. \hfill
(3pts)
\end{enumerate}
\noindent\hfill\textbf{Présentation : 1pt}

% ═══════════════════════════════════════════════════════════════
\section{Session 2024 — Corrigé}
\begin{multicols}{2}

\corrige{— Activités numériques, Ex. 1.1}
$(2+\sqrt3)^{2}=4+4\sqrt3+3=7+4\sqrt3$.

\corrige{— Ex. 1.2}
$A=\dfrac{2+\sqrt3}{2-\sqrt3}=\dfrac{(2+\sqrt3)^{2}}{(2-\sqrt3)(2+\sqrt3)}=\dfrac{7+4\sqrt3}{4-3}=7+4\sqrt3$.

\corrige{— Ex. 1.3}
$4\times1{,}732=6{,}928$ et $4\times1{,}733=6{,}932$, donc $13{,}928<A<13{,}932$ : un
encadrement d'ordre $2$ est $13{,}92<A<13{,}94$.

\corrige{— Ex. 2.1}
$B=(x-2)[(x-2)+(x+4)]=(x-2)(2x+2)=2(x-2)(x+1)$.

\corrige{— Ex. 2.2}
Condition d'existence : $x\ne2$ et $x\ne-3$. $C=\dfrac{2x+1}{x+3}$.

\corrige{— Ex. 3.1}
Caractère étudié : la note sur $20$ en mathématiques ; nature quantitative continue
(regroupée en classes).

\corrige{— Ex. 3.2}
Effectif total $=50$. $13+15+10+5=43$, donc l'effectif de $[12\,;16[$ est $50-43=7$.

\corrige{— Ex. 3.3}
Classe modale (effectif maximal $15$) : $[4\,;8[$.

\corrige{— Ex. 3.4}
Avec les centres de classes $2$, $6$, $10$, $14$, $18$ :\\
moyenne $=\dfrac{26+90+100+98+90}{50}$\\
$=\dfrac{404}{50}=8{,}08$.

\corrige{— Activités géométriques, Ex. 1.1}
$\overrightarrow{AB}=(-3-2\,;1-2)=(-5\,;-1)$. \textbf{Vrai}.

\corrige{— Ex. 1.2}
$AC=\sqrt{(4-2)^{2}+(-2-2)^{2}}=\sqrt{4+16}=\sqrt{20}=2\sqrt5\ne3\sqrt5$. \textbf{Faux}.

\corrige{— Ex. 1.3}
Droite de direction $\overrightarrow{AB}=(-5\,;-1)$ passant par $A(2\,;2)$ :
$x=5y-8\Rightarrow x-5y+8=0$. \textbf{Vrai}.

\corrige{— Ex. 2.1}
$\cos\widehat{BAC}=\dfrac{AC}{AB}=\dfrac45=0{,}8$, donc
$\widehat{BAC}=\arccos(0{,}8)\approx37^{\circ}$.

\corrige{— Ex. 2.2}
$A$, $B$, $C$, $M$ sont sur le même cercle : les angles inscrits $\widehat{BAC}$ et
$\widehat{BMC}$ interceptent le même arc $BC$, donc ils ont la même mesure.

\corrige{— Ex. 2.3}
$\dfrac{AE}{AB}=\dfrac45$. Par Thalès, $EF=BC\times\dfrac45=3\times\dfrac45=2{,}4$ cm.

\corrige{— Ex. 3.1}
$V=\dfrac13\times AB^{2}\times SM=\dfrac13\times16\times6=32$ cm$^{3}$.

\corrige{— Ex. 3.2}
$k=\dfrac{SN}{SM}=\dfrac12$, donc $v=V\times k^{3}=32\times\dfrac18=4$ cm$^{3}$.

\corrige{— Partie B, Tâche 1}
Soit $B$ le budget total. Part banque $+$ part recettes $=\dfrac25B+\dfrac14B=\dfrac{13}{20}B$,
donc le reste $=\dfrac{7}{20}B=1\,225\,000\Rightarrow B=1\,225\,000\times\dfrac{20}7=3\,500\,000$
FCFA.

\corrige{— Tâche 2}
Dimensions en cm : $480\times300$. Plus grand carreau carré possible : côté
$=\text{PGCD}(480\,;300)=60$ cm. Nombre de carreaux $=\dfrac{480}{60}\times\dfrac{300}{60}=8\times5=40$.

\corrige{— Tâche 3}
Soit $x$ le nombre de seaux d'eau et $y$ le nombre de seaux d'huile : $x+y=10$ et
$40\,000x+60\,000y=440\,000$, soit $2x+3y=22$. En substituant $x=10-y$ :
$20-2y+3y=22\Rightarrow y=2$, donc $x=8$. Achat : $8$ seaux de peinture à eau et $2$ seaux
de peinture à huile.

\end{multicols}

% ═══════════════════════════════════════════════════════════════
\section{Session 2025 — Épreuve}

\noindent\textbf{Examen :} BEPC \quad\textbf{Session :} 2025 \quad\textbf{Épreuve :}
Mathématiques \quad\textbf{Durée :} 2 heures \quad\textbf{Coefficient :} 4.

\subsection*{Partie A : évaluation des ressources (12,5 points)}

\paragraph{Activités numériques (6,25 points)}

\textbf{Exercice 1} (3,25 points)
\begin{enumerate}[label=\arabic*)]
\item Écrire le nombre $A=\left(\dfrac32+\dfrac54\right)\times\left(\dfrac32-\dfrac54\right)$
sous la forme d'une fraction irréductible. \hfill (1pt)
\item[2)] a) Écrire le nombre $B=(3\sqrt2-4)^{2}$ sous la forme $a+b\sqrt2$ où $a$ et $b$
sont des entiers. \hfill (0,75pt) \\
b) Montrer que le nombre $C=(3\sqrt2-4)(3\sqrt2+4)$ est un entier. \hfill (0,75pt) \\
c) Écrire le nombre $D=\dfrac{3\sqrt2-4}{3\sqrt2+4}$ sans radical au dénominateur. \hfill
(0,75pt)
\end{enumerate}

\textbf{Exercice 2} (3 points) \\
On considère l'expression $E=(x-3)(x+4)+x^{2}-16$.
\begin{enumerate}[label=\arabic*)]
\item Choisir la forme développée de $E$ parmi celles qui sont proposées ci-dessous. \\
a) $2x^{2}-x-28$~; \quad b) $2x^{2}-x+28$~; \quad c) $2x^{2}+x+28$~; \quad d) $2x^{2}+x-28$.
\hfill (1pt)
\item Choisir la forme factorisée de $E$ parmi celles qui sont proposées ci-dessous : \\
a) $(x-4)(2x-7)$~; \quad b) $(x+4)(2x-7)$~; \quad c) $(x-4)(2x+7)$~; \quad d)
$(x+4)(2x+7)$. \hfill (1pt)
\item[3)] On considère la fraction rationnelle $Q=\dfrac{(x+4)(2x-7)}{(x+4)(2x+7)}$. \\
a) Écrire la condition d'existence d'une valeur numérique de $Q$. \hfill (0,5pt) \\
b) Simplifier $Q$. \hfill (0,5pt)
\end{enumerate}

\paragraph{Activités géométriques (6,25 points)}

\textbf{Exercice 1} (2 points) \\
Répondre par vrai ou faux.
\begin{enumerate}[label=\arabic*)]
\item Dans un repère orthonormé du plan, les droites $(D)$ et $(L)$ d'équations
respectives $y=2x-1$ et $y=\dfrac12x+3$ sont perpendiculaires. \hfill (1pt)
\item Les vecteurs $\vec u$ et $\vec v$ de coordonnées respectives $\binom{-2}{5}$ et
$\binom52$ sont colinéaires. \hfill (1pt)
\end{enumerate}

\textbf{Exercice 2} (2,25 points) \\
$ABC$ est un triangle rectangle en $B$ tel que $AB=6$ cm et $BC=8$ cm.

\begin{illus}
\begin{tikzpicture}[scale=0.55]
\draw[onbo,line width=1.3pt] (0,0) -- (6,0) -- (0,8) -- cycle;
\draw[onbmuted,line width=0.8pt] (0.5,0)--(0.5,0.5)--(0,0.5);
\node[below,font=\small] at (3,0) {$6$ cm};
\node[left,font=\small] at (0,4) {$8$ cm};
\node[onbo,above right,font=\small] at (3,4) {$AC$};
\node[below left,font=\small] at (0,0) {$B$};
\node[below,font=\small] at (6,0) {$C$};
\node[above,font=\small] at (0,8) {$A$};
\end{tikzpicture}
\fig{Figure — Triangle $ABC$ rectangle en $B$, $AB=6$ cm, $BC=8$ cm.}
\end{illus}

\begin{enumerate}[label=\arabic*)]
\item Calculer $AC$. \hfill (1,25pt)
\item Calculer la tangente de l'angle $\widehat{BAC}$ et en déduire la mesure de l'angle
$\widehat{BAC}$ arrondie à l'entier le plus proche. \hfill (1pt)
\end{enumerate}

\textbf{Exercice 3} (2 points) \\
$SABCD$ est une pyramide régulière de hauteur $9$ cm, dont la base est un carré de côté
$8$ cm.

\begin{illus}
\begin{tikzpicture}[scale=0.85]
\coordinate (O) at (0,0);
\coordinate (F) at (0,-1.0);
\coordinate (K) at (0,1.0);
\coordinate (R) at (1.9,-0.35);
\coordinate (L) at (-1.9,0.35);
\coordinate (S) at (0,3.0);
\draw[onbline,dashed,line width=0.7pt] (R) -- (K) -- (L);
\draw[onbline,line width=1pt] (F) -- (R);
\draw[onbline,line width=1pt] (L) -- (F);
\draw[onbline,dashed,line width=0.7pt] (S) -- (K);
\draw[onbo,line width=1.2pt] (S) -- (F);
\draw[onbo,line width=1.2pt] (S) -- (R);
\draw[onbo,line width=1.2pt] (S) -- (L);
\draw[onbmuted,dashed,line width=0.7pt] (S) -- (O);
\node[onbmuted,left,font=\scriptsize] at (0,1.6) {$9$ cm};
\node[below right,font=\scriptsize] at (0.9,-0.8) {$8$ cm};
\node[above,font=\scriptsize] at (S) {$S$};
\end{tikzpicture}
\fig{Figure — Pyramide régulière $SABCD$, hauteur $9$ cm, base carrée de côté $8$ cm.}
\end{illus}

\noindent Calculer le volume de cette pyramide. \hfill (2pts)

\subsection*{Partie B : évaluation des compétences (7,5 points)}

\textbf{Situation.} Une association décide de venir en aide, en fin d'année, à un
orphelinat qui a présenté un budget.

\textbf{I.} En début d'année, la caisse de l'association disposait encore d'une somme de
$500\,000$ FCFA. Les membres de l'association ont cotisé en plus une somme représentant
$\dfrac45$ du budget. Mais au moment de faire les comptes, ils constatent qu'il leur manque
$100\,000$ FCFA pour boucler le budget.

\textbf{Tâche 1 :} Quel est le montant du budget présenté par l'orphelinat ? \hfill
(2,25pts)

\textbf{II.} Les dons en nature constitués de sacs de riz et de cartons d'huile ont été
achetés en deux temps chez le même vendeur :
\begin{itemize}
\item $1^{\text{er}}$ temps : $3$ sacs de riz et $2$ cartons d'huile pour un montant total
de $146\,000$ FCFA ;
\item $2^{\text{ème}}$ temps : $2$ sacs de riz et $2$ cartons d'huile pour un montant total
de $111\,000$ FCFA.
\end{itemize}
Les factures des achats ont été malheureusement égarées.

\textbf{Tâche 2 :} Combien a coûté un sac de riz et un carton d'huile ? \hfill (2,25pts)

\textbf{III.} Le chauffeur du véhicule chargé du transport des achats a proposé deux
options :
\begin{itemize}
\item $1^{\text{ère}}$ option : payer $5\,000$ FCFA pour le carburant et $2\,000$ FCFA par
heure ;
\item $2^{\text{ème}}$ option : payer $10\,000$ FCFA pour le carburant et $1\,000$ FCFA par
heure.
\end{itemize}
Le membre chargé des achats voudrait savoir le nombre d'heures pour lequel la
$1^{\text{ère}}$ option est la plus avantageuse.

\textbf{Tâche 3 :} À partir de quelle durée en heures, l'option $2$ est-elle plus
avantageuse pour le transport des achats ? \hfill (2,25pts)

\noindent\hfill\textbf{Présentation : 0,75pt}

% ═══════════════════════════════════════════════════════════════
\section{Session 2025 — Corrigé}
\begin{multicols}{2}

\corrige{— Activités numériques, Ex. 1.1}
$A=\left(\dfrac32+\dfrac54\right)\left(\dfrac32-\dfrac54\right)=\left(\dfrac32\right)^{2}-\left(\dfrac54\right)^{2}=\dfrac94-\dfrac{25}{16}=\dfrac{36-25}{16}=\dfrac{11}{16}$.

\corrige{— Ex. 1.2a}
$B=(3\sqrt2)^{2}-2\times3\sqrt2\times4+16=18-24\sqrt2+16=34-24\sqrt2$.

\corrige{— Ex. 1.2b}
$C=(3\sqrt2)^{2}-4^{2}=18-16=2$ : entier.

\corrige{— Ex. 1.2c}
$D=\dfrac{(3\sqrt2-4)^{2}}{(3\sqrt2+4)(3\sqrt2-4)}=\dfrac{34-24\sqrt2}{2}=17-12\sqrt2$.

\corrige{— Ex. 2.1}
$E=x^{2}+x-12+x^{2}-16=2x^{2}+x-28$. Réponse : \textbf{d}.

\corrige{— Ex. 2.2}
$E=2x^{2}+x-28=(x+4)(2x-7)$ (on vérifie : $2x^{2}-7x+8x-28=2x^{2}+x-28$). Réponse :
\textbf{b}.

\corrige{— Ex. 2.3}
a) Condition d'existence : $x\ne-4$ et $x\ne-\dfrac72$. \\
b) $Q=\dfrac{2x-7}{2x+7}$.

\corrige{— Activités géométriques, Ex. 1.1}
$(D)$ a pour coefficient directeur $2$, $(L)$ a pour coefficient directeur $\dfrac12$.
Produit $=2\times\dfrac12=1\ne-1$ : les droites ne sont \textbf{pas} perpendiculaires.
\textbf{Faux}.

\corrige{— Ex. 1.2}
$\vec u=\binom{-2}5$, $\vec v=\binom52$ : déterminant $=(-2)\times2-5\times5=-4-25=-29\ne0$.
Les vecteurs ne sont \textbf{pas} colinéaires. \textbf{Faux}.

\corrige{— Ex. 2.1}
$AC=\sqrt{AB^{2}+BC^{2}}=\sqrt{36+64}=\sqrt{100}=10$ cm.

\corrige{— Ex. 2.2}
$\tan\widehat{BAC}=\dfrac{BC}{AB}=\dfrac86\approx1{,}33$, donc
$\widehat{BAC}=\arctan\left(\dfrac86\right)\approx53^{\circ}$.

\corrige{— Ex. 3}
$V=\dfrac13\times8^{2}\times9=\dfrac13\times64\times9=192$ cm$^{3}$.

\corrige{— Partie B, Tâche 1}
Soit $B$ le budget. $500\,000+\dfrac45B+100\,000=B\Rightarrow600\,000=\dfrac15B\Rightarrow
B=3\,000\,000$ FCFA.

\corrige{— Tâche 2}
Soit $r$ le prix d'un sac de riz et $h$ celui d'un carton d'huile : $3r+2h=146\,000$ et
$2r+2h=111\,000$. En soustrayant : $r=35\,000$ FCFA, puis
$2h=111\,000-2\times35\,000=41\,000\Rightarrow h=20\,500$ FCFA.

\corrige{— Tâche 3}
Coût option 1 $=5\,000+2\,000t$ ; coût option 2 $=10\,000+1\,000t$. Option 2 avantageuse si
$10\,000+1\,000t<5\,000+2\,000t\Leftrightarrow5\,000<1\,000t\Leftrightarrow t>5$. À partir de
$6$ heures (première valeur entière supérieure à $5$), l'option $2$ est la plus
avantageuse.

\end{multicols}
\exercicesBanner
